\documentclass{article}

\usepackage{arxiv}

\usepackage[utf8]{inputenc} % allow utf-8 input
\usepackage[T1]{fontenc}    % use 8-bit T1 fonts
\usepackage{hyperref}       % hyperlinks
\usepackage{url}            % simple URL typesetting
\usepackage{booktabs}       % professional-quality tables
\usepackage{amsfonts}       % blackboard math symbols
\usepackage{amsmath}        % align, cases, etc.
\usepackage{nicefrac}       % compact symbols for 1/2, etc.
\usepackage{microtype}      % microtypography
\usepackage{graphicx}
\usepackage{natbib}
\usepackage{doi}

% =========================================================================
%  STRUCTURE-ONLY SCAFFOLD (section titles + float placement notes).
%  No prose yet, by design — we agreed to lock headings first, then draft
%  section-by-section. Every "% [TODO ...]" marks content to be written and
%  every "% [FIG ...]" / "% [TAB ...]" marks a float to place in that section.
% =========================================================================

\title{When Can Scientific Machine Learning Recover a Hidden Ionic Current?
Identifiability Limits of Universal Differential Equations in a
Hodgkin--Huxley Neuron}

% ---- Author block ------------------------------------------------------
% [TODO authors] Confirm exact names, order, affiliations, emails, ORCIDs.
% Decided: Vinayak (first author) + Vizuara team (Prathamesh + other mentors).
\author{
  Vinayak Mokashi \\
  Vizuara \\ % [TODO affiliation/address]
  \texttt{vinayakmokashi305@gmail.com} \\
  \And
  Prathamesh \\ % [TODO full name, affiliation, email, ORCID]
  Vizuara \\
  \texttt{} \\
  %% \And  % [TODO additional Vizuara mentor(s) as needed]
  %% Coauthor \\ Vizuara \\ \texttt{} \\
}

\renewcommand{\shorttitle}{Identifiability Limits of UDEs in a Hodgkin--Huxley Neuron}

\hypersetup{
  pdftitle={When Can Scientific Machine Learning Recover a Hidden Ionic Current? Identifiability Limits of Universal Differential Equations in a Hodgkin-Huxley Neuron},
  pdfsubject={q-bio.NC, q-bio.QM, cs.LG},
  pdfauthor={Vinayak Mokashi},
  pdfkeywords={scientific machine learning, universal differential equations, neural ODE, Hodgkin-Huxley, symbolic regression, SINDy, identifiability, computational neuroscience},
}

\begin{document}
\maketitle

\begin{abstract}
% [TODO abstract] Write last; ~150-200 words. Arc: advanced HH neuron with a
% hidden calcium current -> NODE forecasts fail, physics-embedded UDE recovers
% the trajectory and the hidden current (multi-seed, incl. voltage-only) ->
% symbolic recovery gives interpretable parameters BUT is identifiability-
% limited -> two controls (gCa sweep + aggressive-retrain) show the residual
% gap is a genuine single-trajectory identifiability limit, not a training
% artifact. Headline takeaway = when SciML can (and cannot) recover mechanism.
\end{abstract}

\keywords{Scientific machine learning \and Universal differential equations \and
Neural ODEs \and Hodgkin--Huxley \and Symbolic regression / SINDy \and
Identifiability \and Computational neuroscience}

% =========================================================================
\section{Introduction}
\label{sec:intro}

\subsection{Motivation: partially-known dynamics in conductance-based models}

\subsection{The gap: from prediction to mechanism, and when it fails}

\subsection{Contributions}
% [TODO] 4-item bulleted list: (1) reproducible advanced-HH SciML benchmark
% (NODE vs UDE); (2) physics-embedded UDE recovering the hidden current, robust
% across seeds and under voltage-only observation; (3) symbolic recovery via a
% biophysical conductance form + unconstrained SINDy; (4) identifiability
% analysis with two controls (gCa sweep + aggressive-retrain).

% =========================================================================
\section{Related Work}
\label{sec:related}
% [TODO] Neural ODEs; Universal DEs (Rackauckas et al.); SINDy (Brunton et al.);
% data-driven / physics-informed fitting of conductance-based models;
% parameter identifiability in neuroscience.

% =========================================================================
\section{Background and Problem Setup}
\label{sec:setup}

\subsection{The advanced Hodgkin--Huxley model}
% [FIG] dynamics overview (figs 01-08, composite): V + gates + the two extra currents.

\subsection{Known physics versus the hidden current: the recovery task}

\subsection{Data generation and evaluation protocol}
% [TODO] OU noise; train-on-noisy / evaluate-vs-clean; train/forecast split;
% common evaluation window.

% =========================================================================
\section{Methods}
\label{sec:methods}

\subsection{Black-box Neural ODE baseline}

\subsection{Universal Differential Equation with embedded physics}
% [TODO] full-state + voltage-only (partial-observation) variant.

\subsection{Symbolic recovery of the learned current}
% [TODO] hull-restricted grid fit; unconstrained SINDy; conductance form s^2(aV+b).

\subsection{Training, metrics, and reproducibility}
% [TODO] multi-seed; common-eval window; representative-seed guardrail.

% =========================================================================
\section{Experiments and Results}
\label{sec:results}

\subsection{Forecasting: the Neural ODE diverges, the UDE extrapolates}
% [FIG] fig2 (NODE), fig3 (UDE), fig5 (NODE vs UDE bars, mean+/-std).
% [TAB] headline metrics, mean+/-std over 5 seeds (metrics_baseline_multiseed.csv).

\subsection{Partial observation: the voltage-only UDE}
% [FIG] fig6, fig8 (full vs voltage-only, mean+/-std).

\subsection{Ablations: noise, training window, and calcium identifiability}
% [FIG] fig7 (noise), fig7b/fig7b_commoneval (window), fig7c (gCa).
% [TODO] note the 40-50 ms window non-monotonicity as an optimization artifact.

\subsection{Symbolic recovery of the calcium current}
% [FIG] fig9 (parity), fig10 (timeseries), fig11 (coeff recovery), fig12 (identifiability).
% [TAB] recovered coefficients a, b, E_rev, R^2 (symbolic_recovery_metrics.csv).

\subsection{Is the residual gap trainable? An aggressive-retrain control}
% [TAB] before/after (5k/300 vs 20k/1000): ensemble a, per-seed a std, cond R^2,
% SINDy term count, forecast V-RMSE + the seed-1111 decoupling. => genuine
% single-trajectory identifiability limit, not a training-budget artifact.

% =========================================================================
\section{Discussion}
\label{sec:discussion}
% [TODO] identifiability as the central lesson; success vs. limits of mechanism
% recovery; gCa=2.0 as a controlled experiment (not convenience); window
% non-monotonicity caveat; limitations; when to expect recovery to succeed.

% =========================================================================
\section{Conclusion}
\label{sec:conclusion}

% =========================================================================
\section*{Acknowledgements}
% [TODO] Vizuara mentors (Prathamesh; SciML guidance), compute. NeurIPS-style:
% acknowledgements are unnumbered and typically after the conclusion.

% =========================================================================
\bibliographystyle{unsrtnat}
\bibliography{references} % [TODO] build references.bib

% =========================================================================
\appendix
\section{Model equations and gating kinetics}
% [TODO] full alpha/beta kinetics, p_inf/s_inf, constants + literature refs.

\section{Full hyperparameters and training details}
% [TODO] optimizer schedule, solver/adjoint, seeds, iteration budgets.

\section{Additional results}
% [TODO] per-seed tables; extra figures; reproducibility (repo + seeds).

\end{document}
